\documentclass[../../main.tex]{subfiles}% 注意这里的文件路径不能用 ./main.tex ,否则用latexmk编译子文件会报错
\graphicspath{{\subfix{./image/}}} % 指定图片目录,后续可以直接使用图片文件名
% 注意这里的文件路径不能用 ../../image/ ,否则用latexmk编译子文件会报错

% 例如:
% \begin{figure}[H]
% \centering
% \includegraphics[scale=0.3]{图.png}
% \caption{}
% \label{figure:图}
% \end{figure}
% 注意:上述\label{}一定要放在\caption{}之后,否则引用图片序号会只会显示??.

\begin{document}

\section{$B_0$空间}

\begin{definition}
设$\mathscr{X}$是一个线性空间，称它是\textbf{可数范数空间}(或$B_0^*$空间)，是指在它上面有可数个半范数$\{\|\cdot\|_m\}_{m=1}^\infty$，满足：
\begin{enumerate}[(1)]
\item $\|x+y\|_m\leqslant\|x\|_m+\|y\|_m\quad(\forall x,y\in\mathscr{X})$;
\item $\|\lambda x\|_m=|\lambda|\|x\|_m\quad(\lambda\in\mathbb{R},x\in\mathscr{X})$;
\item $\|x\|_m\geqslant0,\|\theta\|_m=0\quad(\forall x\in\mathscr{X})$;
\item $\|x\|_m=0\ (m=1,2,\cdots)\iff x=\theta$.
\end{enumerate}
\end{definition}
\begin{remark}
实际上每个$\|\cdot\|_m$是半范数。
\end{remark}
\begin{remark}
在可数范数空间定义中，可数个半范数可以换成满足下列条件的可数个半范数：
\begin{align*}
\|x\|_1'\leqslant\|x\|_2'\leqslant\cdots\leqslant\|x\|_m'\leqslant\cdots\quad(\forall x\in\mathscr{X}).
\end{align*}
事实上，只需令
\begin{align*}
\|x\|_m'=\max(\|x\|_1,\cdots,\|x\|_m)\quad(\forall x\in\mathscr{X}).
\end{align*}
\end{remark}

\begin{definition}
在线性空间$\mathscr{X}$上给定两组可数个半范数$\{\|\cdot\|_m\}_{m=1}^\infty$与$\{\|\cdot\|_m'\}_{m=1}^\infty$，如果它们导出相同的收敛性，则称它们是\textbf{等价的}。
\end{definition}

\begin{proposition}
在线性空间$\mathscr{X}$上，为了两组可数个半范数$\{\|\cdot\|_m\}_{m=1}^\infty$与$\{\|\cdot\|_m'\}_{m=1}^\infty$是等价的，必须且仅须：
$\forall m\in\mathbb{N},\exists m'\in\mathbb{N}$及$\exists C_{mm'}>0$，使得
\begin{align*}
\|x\|_m\leqslant C_{mm'}\|x\|_{m'}'\quad(\forall x\in\mathscr{X}),
\end{align*}
并且$\forall n'\in\mathbb{N},\exists n\in\mathbb{N}$及$\exists C_{n'n}'>0$，使得
\begin{align*}
\|x\|_{n'}'\leqslant C_{n'n}'\|x\|_n\quad(\forall x\in\mathscr{X}).
\end{align*}
\end{proposition}

\begin{proposition}
每个$B_0^*$空间$\mathscr{X}$必是一个$F^*$空间，即若$\{\|\cdot\|_m\}_{m=1}^\infty$是可数个半范数，则
\begin{align*}
\|x\|=\sum\limits_{m=1}^\infty\frac{1}{2^m}\cdot\frac{\|x\|_m}{1+\|x\|_m}\quad(\forall x\in\mathscr{X})
\end{align*}
是一个准范数，并且$\|\cdot\|$导出的收敛性与$\{\|\cdot\|_m\}_{m=1}^\infty$导出的收敛性一致。
\end{proposition}

\begin{example}
$\mathscr{D}_K$是$B_0^*$空间，其可数范数$\|\cdot\|_m$按对应定义式规定。
\end{example}

\begin{example}\label{example:泛函分析--例4.2.6}
设$\Omega$是$\mathbb{R}^n$中的任意开集，$K_m$是一串相对于$\Omega$的紧集，满足
\begin{align*}
K_1\subset\mathring{K}_2\subset K_2\subset K_3\subset\cdots\subset K_m\subset\cdots\subset\Omega,\quad \bigcup\limits_{m=1}^\infty K_m=\Omega,
\end{align*}
令
\begin{align*}
\|\varphi\|_m=\sum\limits_{|\alpha|\leqslant m}\max\limits_{x\in K_m}|\partial^\alpha\varphi(x)|\quad(m=1,2,\cdots).
\end{align*}
用$\mathscr{E}(\Omega)$表示带有可数范数$\{\|\cdot\|_m\}_{m=1}^\infty$的线性空间$C^\infty(\Omega)$，则$\mathscr{E}(\Omega)$是一个$B_0^*$空间。$\varphi_j\to0(\mathscr{E}(\Omega))$当且仅当$\forall\varepsilon>0,\forall m\in\mathbb{N},\exists N=N(\varepsilon,m)$，使得
\begin{align*}
\max\limits_{x\in K_m}|\partial^\alpha\varphi_j(x)|<\varepsilon\quad(\forall|\alpha|\leqslant m,\forall j>N),
\end{align*}
即在每一个相对于$\Omega$的紧集$K_m$上，直到$m$次导数一致收敛于$0$。

$\mathscr{E}(\Omega)$上的收敛性与紧集列$\{K_m\}$的选择无关，即按不同的紧集列定义出的两组可数个半范数是等价的。
\end{example}

\begin{example}
用$\mathscr{S}(\mathbb{R}^n)$表示集合
\begin{align*}
\{\varphi\in C^\infty(\mathbb{R}^n)\mid\sup\limits_{x\in\mathbb{R}^n}(1+|x|^2)^{\frac{k}{2}}|\partial^\alpha\varphi(x)|\leqslant M_{k,\alpha}<\infty\quad(k,|\alpha|=0,1,2,\cdots)\}.
\end{align*}
定义半范数为
\begin{align*}
\|\varphi\|_m=\sup\limits_{\substack{|\alpha|\leqslant m\\x\in\mathbb{R}^n}}\left|(1+|x|^2)^{\frac{m}{2}}\partial^\alpha\varphi(x)\right|\quad(m=0,1,2,\cdots).
\end{align*}
$\mathscr{S}(\mathbb{R}^n)$上的函数称为\textbf{速降函数}，其任意阶导数在无穷远处比任何负幂次下降得都快，即对任意非负整数$m$及多重指标$\alpha$，都有
\begin{align*}
\lim\limits_{|x|\to\infty}(1+|x|^2)^{\frac{m}{2}}|\partial^\alpha\varphi(x)|=0.
\end{align*}
\end{example}

\begin{definition}
一个$B_0^*$空间$\mathscr{X}$称为是\textbf{完备的}，是指其中的任何基本列都是收敛的。完备的$B_0^*$空间称为$B_0$空间。
\end{definition}

\begin{proposition}
$\mathscr{S}(\mathbb{R}^n)$是$B_0$空间。
\end{proposition}
\begin{proof}
设$\{\varphi_\nu(x)\}$是$\mathscr{S}(\mathbb{R}^n)$中的一个基本列。由定义，对$\forall m\in\mathbb{N}$及$\forall\varepsilon>0,\exists N=N(m,\varepsilon)$，当$\mu,\nu>N$时，
\begin{align}
\sup\limits_{\substack{|\alpha|\leqslant m\\x\in\mathbb{R}^n}}\left|(1+|x|^2)^{\frac{m}{2}}\partial^\alpha(\varphi_\mu-\varphi_\nu)(x)\right|<\varepsilon.\label{eq:4-2-2}
\end{align}
\begin{enumerate}[(1)]
\item 因对每个$m$，$\{\|\varphi_\nu\|_m\}$是有界的，故必存在常数$M_m$，使得
\begin{align*}
\sup\limits_{\substack{|\alpha|\leqslant m\\x\in\mathbb{R}^n}}\left|(1+|x|^2)^{\frac{m}{2}}\partial^\alpha\varphi_\nu(x)\right|\leqslant M_m\quad(\nu=1,2,\cdots),
\end{align*}
从而
\begin{align}
|\partial^\alpha\varphi_\nu(x)|\leqslant\frac{M_m}{(1+|x|^2)^{m/2}}.\label{eq:4-2-3}
\end{align}
此外，在任意有界闭球$|x|\leqslant R$上，$\{\partial^\alpha\varphi_\nu(x)\}$依一致范数是基本序列，故在$|x|\leqslant R$上，$\{\partial^\alpha\varphi_\nu(x)\}$有一个一致极限$\psi_\alpha(x)$，满足
\begin{align}
|\psi_\alpha(x)|\leqslant\frac{M_m}{(1+|x|^2)^{m/2}}\quad(|\alpha|\leqslant m).\label{eq:4-2-4}
\end{align}

\item 可证$\psi_\alpha(x)=\partial^\alpha\psi_0(x)$，只需证明$\psi_{(1,0,\cdots,0)}(x)=\partial_{x_1}\psi_0(x)$，其余用归纳法递推。
对任意的$R>0$，当$|x|\leqslant R$时，有
\begin{align*}
\partial_{x_1}\varphi_\nu(x)\rightrightarrows\psi_{(1,0,\cdots,0)}(x)\quad(\nu\to\infty),
\end{align*}
以及
\begin{align*}
\varphi_\nu(x)\rightrightarrows\psi_0(x)\quad(\nu\to\infty).
\end{align*}
对$\forall\varepsilon>0$，取$M_1'>0$及$N_0\in\mathbb{N}$，使得
\begin{align*}
\int_{|x_1|>M_1'}\frac{\mathrm{d}x_1}{(1+|x_1|^2)^{1/2}}<\frac{\varepsilon}{4M_1'},
\end{align*}
以及
\begin{align*}
\left|\psi_{(1,0,\cdots,0)}(x)-\partial_{x_1}\varphi_\nu(x)\right|<\frac{\varepsilon}{4M_1'}\quad(|x|\leqslant M_1',\nu>N_0),
\end{align*}
便有
\begin{align*}
&\left|\int_{-\infty}^{x_1}\psi_{(1,0,\cdots,0)}(x',x_2,\cdots,x_n)\mathrm{d}x'-\varphi_\nu(x_1,x_2,\cdots,x_n)\right|\\
=&\left|\int_{-\infty}^{x_1}\left[\psi_{(1,0,\cdots,0)}(x',x_2,\cdots,x_n)-\partial_{x_1}\varphi_\nu(x',x_2,\cdots,x_n)\right]\mathrm{d}x'\right|\\
\leqslant&\int_{|x'|>M_1'}\left|\psi_{(1,0,\cdots,0)}(x',x_2,\cdots,x_n)\right|\mathrm{d}x'+\int_{|x'|>M_1'}\left|\partial_{x_1}\varphi_\nu(x',x_2,\cdots,x_n)\right|\mathrm{d}x'\\
&+\int_{|x'|\leqslant M_1'}\left|\psi_{(1,0,\cdots,0)}(x',x_2,\cdots,x_n)-\partial_{x_1}\varphi_\nu(x',x_2,\cdots,x_n)\right|\mathrm{d}x'\\
<&2M_1\cdot\frac{\varepsilon}{4M_1}+2M_1'\cdot\frac{\varepsilon}{4M_1'}=\varepsilon.
\end{align*}
从而
\begin{align*}
\psi_0(x)=\int_{-\infty}^{x_1}\psi_{(1,0,\cdots,0)}(x',x_2,\cdots,x_n)\mathrm{d}x',
\end{align*}
即
\begin{align*}
\psi_{(1,0,\cdots,0)}(x)=\partial_{x_1}\psi_0(x).
\end{align*}
同时可证$\psi_0(x)\in\mathscr{S}(\mathbb{R}^n)$。

\item 可证$\|\varphi_\nu-\psi_0\|_m\to0\ (\nu\to\infty)$。
对$\forall\varepsilon>0,\exists R>0$，使得当$|x|>R$时，
\begin{align*}
\sup\limits_{|\alpha|\leqslant m}\left|\partial^\alpha(\varphi_\nu(x)-\psi_0(x))\right|\leqslant\frac{M_m}{(1+|x|^2)^{m/2}}<\varepsilon.
\end{align*}
再确定$N$，使得当$\nu>N$时，在$|x|\leqslant R$上有
\begin{align*}
\sup\limits_{|\alpha|\leqslant m}\left|\partial^\alpha(\varphi_\nu(x)-\psi_0(x))\right|<\varepsilon.
\end{align*}
从而
\begin{align*}
\partial^\alpha\varphi_\nu(x)\rightrightarrows\partial^\alpha\psi_0(x)\quad(\nu\to\infty,\forall x\in\mathbb{R}^n).
\end{align*}
因此，在\eqref{eq:4-2-2}式中令$\mu\to\infty$，即得
\begin{align*}
\|\varphi_\nu-\psi_0\|\leqslant\varepsilon\quad(\forall\nu>N).
\end{align*}
于是空间$\mathscr{S}(\mathbb{R}^n)$是完备的。
\end{enumerate}

\end{proof}

\begin{proposition}
为了$\varphi_\nu\to\varphi_0(\mathscr{D}(\Omega))$，必须且仅须存在紧集$K\subseteq\Omega$，使得$\varphi_0,\varphi_\nu\subset\mathscr{D}_K$，而且
\begin{align*}
\|\varphi_\nu-\varphi_0\|_{m,k}\to0\quad(\nu\to\infty,m=1,2,\cdots),
\end{align*}
其中$\|\varphi\|_{m,k}$是$\mathscr{D}_K$上的可数范数。
\end{proposition}

$\mathscr{D}'(\Omega)$表示广义函数全体，它是$\mathscr{D}(\Omega)$的共轭空间。对于$\mathscr{D}_K(\Omega),\mathscr{E}(\Omega),\mathscr{S}(\mathbb{R}^n)$，对应的共轭空间分别记作$\mathscr{D}_K'(\Omega),\mathscr{E}'(\Omega),\mathscr{S}'(\mathbb{R}^n)$，特别当$\Omega=\mathbb{R}^n$时，简记为$\mathscr{D}_K',\mathscr{E}',\mathscr{S}'$，显然有
\begin{align*}
\mathscr{E}'\subseteq\mathscr{S}'\subseteq\mathscr{D}_K'.
\end{align*}

\begin{lemma}\label{lemma:泛函分析--引理hweiowj}
设$\mathscr{X}$是一个$B_0$空间。为了$\mathscr{X}$上的线性泛函$f$是连续的，必须且仅须$\exists m\in\mathbb{N}$及$M_m>0$，使得
\begin{align*}
|\langle f,\varphi\rangle|\leqslant M_m\|\varphi\|_m\quad(\forall\varphi\in\mathscr{X}).
\end{align*}
\end{lemma}
\begin{proof}
充分性显然。必要性用反证法。倘若不然，对$\forall n\in\mathbb{N},\exists x_n\in\mathscr{X}$，使得
\begin{align}
|\langle f,x_n\rangle|>n\|x_n\|_n.\label{eq:4-2-5}
\end{align}
\begin{enumerate}[(1)]
\item 若$\exists N\in\mathbb{N}$，使得$\|x_n\|_n\neq0\ (\forall n\geqslant N)$，令
\begin{align*}
y_n\triangleq\frac{x_n}{n\|x_n\|_n}\quad(n\geqslant N),
\end{align*}
则
\begin{align*}
\|y_n\|_p=\frac{\|x_n\|_p}{n\|x_n\|_n}\leqslant\frac{1}{n}\quad(\text{当 }n\geqslant\max(N,p)).
\end{align*}
故$y_n\to\theta\ (n\to\infty)$，但$|\langle f,y_n\rangle|>1\ (n\geqslant N)$，这与$f$的连续性矛盾。

\item 若存在$\{n_i\}_{i=1}^\infty$，使得$\|x_{n_i}\|_{n_i}=0\ (i=1,2,\cdots)$，则有$|\langle f,x_{n_i}\rangle|>0$，令
\begin{align*}
y_{n_i}=\frac{x_{n_i}}{\langle f,x_{n_i}\rangle}\quad(i=1,2,\cdots),
\end{align*}
便有$\langle f,y_{n_i}\rangle=1$，但$\|y_{n_i}\|_{n_i}=0\ (i=1,2,\cdots)$，这也与$f$的连续性矛盾。
\end{enumerate}

\end{proof}

\begin{theorem}
任意$B_0$空间$\mathscr{X}$具有足够多的连续线性泛函。
\end{theorem}
\begin{proof}
若$f_0$在$\mathscr{X}$的一个线性闭子空间$\mathscr{X}_0$上有定义且连续线性，则存在半范数$\|\cdot\|_m$及常数$M_m>0$，使得
\begin{align*}
|\langle f_0,x\rangle|\leqslant M_m\|x\|_m\quad(\forall x\in\mathscr{X}_0).
\end{align*}
对于半范数$\|\cdot\|_m$，应用Hahn-Banach定理，$f_0$可以扩张到全空间$\mathscr{X}$，即$f$满足$f|_{\mathscr{X}_0}=f_0$，且
\begin{align*}
|\langle f,x\rangle|\leqslant M_m\|x\|_m\quad(\forall x\in\mathscr{X}),
\end{align*}
从而$f\in\mathscr{X}^*$。于是对$\forall x_0\neq\theta,\exists f\in\mathscr{X}^*$，使得$\langle f,x_0\rangle=1$。
\end{proof}

\begin{proposition}
如果有一个$B_0$空间$\mathscr{X}$，满足$\mathscr{D}(\Omega)\hookrightarrow\mathscr{X}$，即$\mathscr{D}(\Omega)\subseteq\mathscr{X}$，并且
\begin{align*}
\varphi_j\to\varphi_0(\mathscr{D}(\Omega))\implies\varphi_j\to\varphi_0(\mathscr{X})\quad(j\to\infty),
\end{align*}
那么$\mathscr{X}$上的任意一个连续线性泛函$f$，必有$f\in\mathscr{D}'(\Omega)$。
\end{proposition}

\begin{example}
$\mathscr{X}=\mathscr{E}(\Omega)$或$L^p(\Omega)(1\leqslant p<\infty)$或$C^k(\Omega)(k\in\mathbb{N})$等都满足$\mathscr{D}(\Omega)\hookrightarrow\mathscr{X}$，从而有$\mathscr{E}'(\Omega),L^q(\Omega)(\frac{1}{p}+\frac{1}{q}=1),[C^k(\Omega)]^*$等都包含于$\mathscr{D}'(\Omega)$。
\end{example}

\begin{theorem}
为了$f\in\mathscr{S}'$，必须且仅须$\exists m\in\mathbb{N}$及$u_\alpha\in L^2(\mathbb{R}^n)(|\alpha|\leqslant m)$，使得
\begin{align*}
\langle f,\varphi\rangle=\sum\limits_{|\alpha|\leqslant m}\int_{\mathbb{R}^n}u_\alpha(x)\partial^\alpha\varphi(x)(1+|x|^2)^{\frac{m}{2}}\mathrm{d}x\quad(\forall\varphi\in\mathscr{S}).
\end{align*}
\end{theorem}
\begin{proof}
充分性显然，只需证必要性。
\begin{enumerate}[(1)]
\item 先证明
\begin{align*}
\|\varphi\|_m'\triangleq\left(\sum\limits_{|\alpha|\leqslant m}\int_{\mathbb{R}^n}(1+|x|^2)^m|\partial^\alpha\varphi(x)|^2\mathrm{d}x\right)^{\frac{1}{2}}\quad(m=1,2,\cdots)
\end{align*}
是$\mathscr{S}$上的一组等价范数。这是因为
\begin{align*}
\|\varphi\|_m'^2&\leqslant\sum\limits_{|\alpha|\leqslant m}\sup\limits_{x\in\mathbb{R}^n}(1+|x|^2)^{m+n}|\partial^\alpha\varphi(x)|^2\int_{\mathbb{R}^n}\frac{\mathrm{d}x}{(1+|x|^2)^n}\\
&\leqslant c\|\varphi\|_{m+n}^2,
\end{align*}
以及
\begin{align*}
\|\varphi\|_m&=\sup\limits_{\substack{|\alpha|\leqslant m\\x\in\mathbb{R}^n}}\left|(1+|x|^2)^{\frac{m}{2}}\partial^\alpha\varphi(x)\right|\\
&\leqslant\sup\limits_{|\alpha|\leqslant m}\int_{\mathbb{R}^n}\left|\frac{\partial^n}{\partial x_1\cdots\partial x_n}(1+|x|^2)^{\frac{m}{2}}\partial^\alpha\varphi(x)\right|\mathrm{d}x\\
&\leqslant c\sup\limits_{|\alpha|\leqslant m}\int_{\mathbb{R}^n}(1+|x|^2)^{\frac{m}{2}}|\partial^{\alpha+e}\varphi(x)|\mathrm{d}x\quad(\text{其中 }e=(1,\cdots,1))\\
&\leqslant C\sum\limits_{|\alpha|\leqslant m+n}\left(\int_{\mathbb{R}^n}(1+|x|^2)^{m+n}|\partial^\alpha\varphi(x)|^2\mathrm{d}x\right)^{\frac{1}{2}}\times\left(\int_{\mathbb{R}^n}\frac{\mathrm{d}x}{(1+|x|^2)^n}\right)^{\frac{1}{2}}\\
&\leqslant c_1\|\varphi\|_{m+n}'.
\end{align*}

\item 对$f$应用\reflem{lemma:泛函分析--引理hweiowj}，$\exists m\in\mathbb{N}$，使得
\begin{align}
|\langle f,\varphi\rangle|\leqslant c_m\|\varphi\|_m'\quad(\forall\varphi\in\mathscr{S}).\label{eq:4-2-6}
\end{align}
于是$f$可以连续地扩张到以$\|\cdot\|_m'$为范数的Banach空间$\mathscr{X}_m$上。注意到$\|\cdot\|_m'$满足平行四边形等式，故可引出内积
\begin{align*}
(\varphi,\psi)_m=\sum\limits_{|\alpha|\leqslant m}\int_{\mathbb{R}^n}\partial^\alpha\varphi(x)\cdot\partial^\alpha\psi(x)(1+|x|^2)^{\frac{m}{2}}\mathrm{d}x,
\end{align*}
该空间构成Hilbert空间。

\item 应用Riesz表示定理，$f$对应着$u\in\mathscr{X}_m$，使得
\begin{align*}
\langle f,\varphi\rangle=\sum\limits_{|\alpha|\leqslant m}\int_{\mathbb{R}^n}\partial^\alpha u(x)\cdot\partial^\alpha\varphi(x)(1+|x|^2)^{\frac{m}{2}}\mathrm{d}x,
\end{align*}
其中
\begin{align*}
\int_{\mathbb{R}^n}|\partial^\alpha u(x)|^2(1+|x|^2)^m\mathrm{d}x<\infty.
\end{align*}
令
\begin{align*}
u_\alpha(x)=(1+|x|^2)^{\frac{m}{2}}\partial^\alpha u(x)\quad(|\alpha|\leqslant m),
\end{align*}
便有$u_\alpha\in L^2(\mathbb{R}^n)$，并且对$\forall\varphi\in\mathscr{S}$，有
\begin{align}
\langle f,\varphi\rangle=\sum\limits_{|\alpha|\leqslant m}\int_{\mathbb{R}^n}u_\alpha(x)\partial^\alpha\varphi(x)(1+|x|^2)^{\frac{m}{2}}\mathrm{d}x.\label{eq:4-2-7}
\end{align}
\end{enumerate}

\end{proof}

\begin{example}
验证: 在\refexa{example:泛函分析--例4.2.6}中，$\mathscr{E}(\Omega)$上的收敛性与紧子集$\{K_m\}$的特殊选择无关。
\end{example}

\begin{example}
设$\|\varphi\|'_m=\sup\limits_{\substack{|k|,|\alpha|\leqslant m\\x\in\mathbb{R}^n}}|x^k\partial^\alpha\varphi(x)|\ (m=0,1,2,\cdots)$，求证: $\|\cdot\|'_m$是$\mathscr{S}(\mathbb{R}^n)$上的等价可数范数。
\end{example}

\begin{example}
验证: $\mathscr{D}_K(\Omega)$与$\mathscr{E}(\Omega)$都是$B_0$空间。
\end{example}

\begin{example}
设$G$是复平面中的有界开连通区域。记$A(G)$为$G$上的解析函数全体，按下列方式规定可数半范数组成的空间: 设
\begin{align*}
G_1\subseteq\overline{G_1}\subseteq G_2\subseteq\overline{G_2}\subseteq\cdots\subseteq G_m\subseteq\overline{G_m}\subseteq\cdots\subseteq G
\end{align*}
是一列连通集，$G_m(m=1,2,\cdots)$是开的，其边界由有穷多段可求长的曲线围成，满足: $\bigcup\limits_{m=1}^\infty\overline{G_m}=G$。令
\begin{align*}
\|\varphi\|_m=\max_{z\in\overline{G_m}}|\varphi(z)| \quad (\forall\varphi\in A(G)).
\end{align*}
求证: $A(G)$是$B_0$空间，又若$\{\varphi_n\}_{n=1}^\infty\subseteq A(G)$，有数列$\{M_m\}_{m=1}^\infty$，使得
\begin{align*}
\|\varphi_n\|_m\leqslant M_m \quad (m=1,2,\cdots;n=1,2,\cdots),
\end{align*}
则$\{\varphi_n\}_{n=1}^\infty$必有收敛子列。
\end{example}




\end{document}